% Generated draft: requires mathematical and source review.
\documentclass[11pt]{article}
\usepackage[margin=25mm]{geometry}
\usepackage{fontspec}
\setmainfont{Latin Modern Roman}
\usepackage{amsmath,amssymb,mathtools,unicode-math}
\setmathfont{Latin Modern Math}
\usepackage{enumitem,needspace,xurl,hyperref}
\setlength{\parindent}{0pt}
\setlength{\parskip}{5pt}
\setlength{\emergencystretch}{3em}
\newcommand{\N}{\mathbb{N}}
\newcommand{\Z}{\mathbb{Z}}
\newcommand{\Q}{\mathbb{Q}}
\newcommand{\R}{\mathbb{R}}
\newcommand{\C}{\mathbb{C}}
\newcommand{\F}{\mathbb{F}}
\newcommand{\A}{\mathbb{A}}
\newcommand{\PP}{\mathbb{P}}
\newcommand{\E}{\mathbb{E}}
\newcommand{\eps}{\varepsilon}
\newcommand{\dd}{\,\mathrm{d}}
\begin{document}
\section*{UnsolvedMath ID: 800366}
\textbf{Source problem code:} catalogue entry
\subsection{Definitions and mathematical statement}
Let \(W\) be a Weyl group of type \(A\), equipped with its standard Coxeter generators and Bruhat order \(\leq\). For \(u,v\in W\) with \(u\leq v\), write
\[
[u,v]=\{x\in W:u\leq x\leq v\}.
\]
Let \(P_{u,v}(q)\) denote the ordinary Kazhdan--Lusztig polynomial, with \(q\) an indeterminate and normalization \(P_{w,w}(q)=1\).

The supplied statement of Brenti's Conjecture 4 asserts that, for every such \(W,u,v\) and every \(z\in[u,v]\) for which \([u,z]\) is a hypercube decomposition of \([u,v]\),
\[
P_{u,v}(q)=\widetilde{P}_{u,v,z}(q).
\]
The equality is an identity in the indeterminate \(q\).

The supplied record does not define either the condition that \([u,z]\) be a hypercube decomposition of \([u,v]\) or the expression \(\widetilde{P}_{u,v,z}(q)\). These source-specific notions remain unresolved here; in particular, the displayed assertion does not apply to every lower interval merely because it has the form \([u,z]\). A self-contained precise formulation requires the definitions and any accompanying hypotheses from the cited source.
\subsection{Short English statement}
Brenti's supplied conjecture says that, in every type \(A\) Weyl group, the Kazhdan--Lusztig polynomial of a Bruhat interval equals the source's expression \(\widetilde{P}_{u,v,z}(q)\) for every hypercube decomposition \([u,z]\) of that interval. The record omits the definitions of hypercube decomposition and of the expression on the right, so the formulation requires source clarification and human mathematical review.
\subsection{Sources}
\begin{itemize}
\item Local UnsolvedMath registry, permanent ID \(800366\), source problem code \texttt{catalogue entry}, title ``Coxeter Groups and Unimodality --- Conjecture 4.'' The supplied record attributes the conjecture to Brenti, with proposed year 2024, and labels it open; this status is supplied metadata, not an independently verified mathematical assessment.
\item Brenti, \emph{Some Open Problems on Coxeter Groups and Unimodality} (2024), Conjecture 4. Supplied source URL: \url{https://arxiv.org/abs/2410.09897}. The provenance records access on 2026-10-16 and extraction by \texttt{arxiv\_latex}; these are supplied provenance claims, not a claim of independent access or verification in preparing this draft.
\item Supplied upstream revision: \texttt{d371ee917ac9647fbb06cd6fc529deec6f3dda82}. The registry record was created and updated on 2026-10-16.
\end{itemize}
\end{document}
